\documentclass[notes,10pt,aspectratio=169]{beamer}
\usepackage{pgfpages}
\setbeameroption{hide notes}
\usepackage{helvet}
\usepackage[default]{lato}
\usepackage{amsmath,amssymb,booktabs,array,graphicx,tikz,pgfplots,eurosym}
\pgfplotsset{compat=1.18}
\useinnertheme[shadow]{rounded}
\definecolor{blue}{RGB}{0,114,178}
\definecolor{red}{RGB}{213,94,0}
\definecolor{yellow}{RGB}{240,228,66}
\definecolor{green}{RGB}{0,158,115}
\setbeamercolor{block title example}{fg=black,bg=orange!50!white}
\setbeamercolor{block body example}{fg=black,bg=orange!30!white}
\setbeamercolor{block title}{use=structure,fg=structure.fg,bg=structure.fg!20!bg}
\setbeamercolor{block body}{parent=normal text,use=block title,bg=block title.bg!50!bg}
\setbeamercolor{section in head/foot}{fg=blue}
\setbeamercolor{section in toc}{fg=blue}
\setbeamercolor{subsection in toc}{fg=red}
\setbeamercolor{frametitle}{fg=blue}
\setbeamercolor{title}{fg=black}
\setbeamertemplate{navigation symbols}{}
\setbeamertemplate{itemize items}{-}
\setbeamercolor{itemize item}{fg=blue}
\setbeamercolor{itemize subitem}{fg=blue}
\setbeamercolor{enumerate item}{fg=blue}
\setbeamercolor{enumerate subitem}{fg=blue}
\setbeamersize{text margin left=0.7cm,text margin right=0.7cm}
\setbeamertemplate{footline}{%
  \hbox{\begin{beamercolorbox}[wd=\paperwidth,ht=3ex,dp=1.5ex,leftskip=2ex,rightskip=2ex]{page footer}%
    \insertsectionnavigationhorizontal{0.7\paperwidth}{}{}\hfill
    \insertframenumber{} / \inserttotalframenumber
  \end{beamercolorbox}}}
\hypersetup{colorlinks=false,linkbordercolor={white},linkcolor={blue}}
\newcommand{\xl}[1]{{\small\texttt{\detokenize{#1}}}}
\newcommand{\datapath}{data/lecture4v3/school_data.csv}
\pgfplotsset{schoolplot/.style={width=11.9cm,height=5.4cm,
  axis lines=left,tick label style={font=\small},label style={font=\small},
  xlabel={Average class size (students)},xmin=15.5,xmax=28.5,
  xtick={16,18,20,22,24,26,28},grid=major,grid style={gray!15}}}
\title[]{\textcolor{blue}{Lecture 4 v3:\\Running a Regression in Excel}}
\author{25117 - Econometrics}
\institute[shortinst]{Universitat Pompeu Fabra}
\date{Academic year 2026--2027}
\begin{document}

\begin{frame}[plain]
\maketitle
\end{frame}

\section{Data}
\begin{frame}{From 100 schools to a regression result}
\begin{enumerate}
\setlength{\itemsep}{9pt}
\item Open and check a dataset with 100 synthetic schools.
\item Draw a scatterplot and estimate a simple regression in Excel.
\item Read coefficients, fit measures, standard errors and confidence intervals.
\item Calculate fitted values and residuals yourself.
\item Introduce family income, $Z$, to prepare for Lecture 5.
\end{enumerate}
\begin{block}{What you need}
Desktop Excel for Windows or Mac and \xl{school_regression.xlsx}. Use the \xl{Data} tab first; \xl{Worked} contains the formula-based answers.
\end{block}
\end{frame}

\begin{frame}{Our question: do schools with smaller classes score higher?}
\begin{itemize}
\setlength{\itemsep}{9pt}
\item One observation is one school, not one pupil: $n=100$ schools.
\item $Y$: the school's average test score, measured in points.
\item $X$: average class size, measured in students per class. A school average can be a decimal.
\item Later, $Z$: average annual family income at the school, measured in \textbf{thousands of euros}.
\end{itemize}
\begin{block}{Interpretation}
These are simulated data designed to resemble an observational school study. A difference in scores or an estimated slope does not, by itself, establish causality.
\end{block}
\end{frame}

\begin{frame}{The dataset: one row per school, one column per variable}
\begin{center}
\begin{tabular}{rrrr}
\toprule
A: school ID & B: score $y_i$ & C: class size $x_i$ & D: income $z_i$\\
\midrule
1 & 610.8186 & 27.8 & 23.39\\
2 & 643.7577 & 19.7 & 33.48\\
3 & 658.5047 & 17.6 & 35.09\\
4 & 663.7172 & 16.3 & 36.68\\
5 & 642.3621 & 24.6 & 27.42\\
6 & 660.9807 & 21.0 & 46.49\\
$\vdots$ & $\vdots$ & $\vdots$ & $\vdots$\\
100 & \multicolumn{3}{c}{The workbook contains all 100 observations.}\\
\bottomrule
\end{tabular}
\end{center}
\begin{itemize}
\item Row 1 contains labels; rows 2--101 contain the observations.
\item Start with columns B and C. Keep column D for the end of the lecture.
\item The numbers are fixed: everyone should obtain the same answers.
\end{itemize}
\end{frame}

\begin{frame}{The model we will estimate}
\[
Y_i=\beta_0+\beta_1X_i+U_i.
\]
\begin{itemize}
\setlength{\itemsep}{8pt}
\item $\beta_0$: intercept; $\beta_1$: slope; $U_i$: factors not included in $X_i$.
\item Excel chooses estimates that minimize the sum of squared residuals:
\[
\min_{b_0,b_1}\sum_{i=1}^{100}(y_i-b_0-b_1x_i)^2.
\]
\item After estimation, $\hat y_i=\hat\beta_0+\hat\beta_1x_i$ and $\hat u_i=y_i-\hat y_i$.
\item We first describe the association between class size and scores. A causal interpretation requires additional assumptions.
\end{itemize}
\end{frame}

\begin{frame}{Step 1. Open the file and check its structure}
\begin{enumerate}
\setlength{\itemsep}{8pt}
\item Open \xl{school_regression.xlsx} in desktop Excel. Save a working copy.
\item Select \xl{Data}. Check the labels in A1:D1 and the last school in row 101.
\item Click B2: the formula bar should show a number, not text or a random formula.
\item Keep each school's values on the same row. If sorting, select the entire table.
\end{enumerate}
\begin{block}{Excel language and regional settings}
The slides use English menu and function names, decimal points, and commas between function arguments. A localized Excel installation may use translated names and semicolons. The supplied workbook avoids CSV import issues.
\end{block}
\end{frame}

\begin{frame}{Step 2. Check the number of observations and basic statistics}
In an empty cell, type a formula beginning with \xl{=}, then press Enter.
\begin{center}
\begin{tabular}{lll}
\toprule
Check & Excel formula & Expected value\\
\midrule
Number of scores & \xl{=COUNT(Data!B2:B101)} & 100\\
Number of class sizes & \xl{=COUNT(Data!C2:C101)} & 100\\
Sample mean of $Y$ & \xl{=AVERAGE(Data!B2:B101)} & 652.9419\\
Sample mean of $X$ & \xl{=AVERAGE(Data!C2:C101)} & 21.7280\\
Sample SD of $Y$ & \xl{=STDEV.S(Data!B2:B101)} & 20.8846\\
Minimum class size & \xl{=MIN(Data!C2:C101)} & 16.1\\
Maximum class size & \xl{=MAX(Data!C2:C101)} & 27.9\\
\bottomrule
\end{tabular}
\end{center}
If a count is below 100, check for missing values or numbers stored as text.
\end{frame}

\section{Excel OLS}
\begin{frame}{Step 3. Create an XY scatterplot}
\begin{enumerate}
\setlength{\itemsep}{7pt}
\item Select \xl{Data!B1:C101}. Choose \textbf{Insert $\rightarrow$ Scatter (X, Y) $\rightarrow$ markers only}.
\item Check the axes: class size must be horizontal; score must be vertical.
\item To set them explicitly: select the chart, then \textbf{Chart Design $\rightarrow$ Select Data}. Edit the series (remove any extra series).
\end{enumerate}
\begin{block}{Series settings}
Series name: \xl{Test score}\\[3pt]
Series X values: \xl{=Data!$C$2:$C$101}\\[3pt]
Series Y values: \xl{=Data!$B$2:$B$101}
\end{block}
Add axis titles: \textbf{Average class size (students)} and \textbf{Test score (points)}.
\end{frame}

\begin{frame}{What the scatterplot should look like}
\centering
\begin{tikzpicture}
\begin{axis}[schoolplot,ylabel={Test score (points)},ymin=595,ymax=710]
\addplot[only marks,mark=*,mark size=1.5pt,blue] table[x=X_class_size,y=Y_score,col sep=comma]{\datapath};
\end{axis}
\end{tikzpicture}
\par\small Each dot is one school. The association is negative, with substantial variation around it.
\end{frame}

\begin{frame}{Step 4. Add a fitted line to the chart}
\begin{columns}[T]
\begin{column}{.43\textwidth}
\begin{enumerate}
\setlength{\itemsep}{8pt}
\item Select a point in the chart.
\item Right-click and choose \textbf{Add Trendline}.
\item Select \textbf{Linear}.
\item Check \textbf{Display Equation on chart} and \textbf{Display R-squared value on chart}.
\end{enumerate}
\vspace{6pt}
Do not force the intercept to zero.
\end{column}
\begin{column}{.54\textwidth}
\begin{block}{Expected equation}
\[
\hat y_i=749.0587-4.4236x_i
\]
\[
R^2=0.5977.
\]
\end{block}
The chart may round its labels. Use regression output or worksheet formulas for calculations.
\vspace{8pt}

A trendline is a quick check. The Regression tool also reports standard errors and confidence intervals.
\end{column}
\end{columns}
\end{frame}

\begin{frame}{Step 5. Enable the Analysis ToolPak}
\begin{columns}[T]
\begin{column}{.48\textwidth}
\begin{block}{Windows}
\begin{enumerate}
\item \textbf{File $\rightarrow$ Options $\rightarrow$ Add-ins}.
\item At the bottom, choose \textbf{Manage: Excel Add-ins $\rightarrow$ Go}.
\item Check \textbf{Analysis ToolPak}; click \textbf{OK}.
\end{enumerate}
\end{block}
\end{column}
\begin{column}{.48\textwidth}
\begin{block}{Mac}
\begin{enumerate}
\item \textbf{Tools $\rightarrow$ Excel Add-ins}.
\item Check \textbf{Analysis ToolPak}; click \textbf{OK}.
\item Restart Excel if necessary.
\end{enumerate}
\end{block}
\end{column}
\end{columns}
\vspace{9pt}
Now look for \textbf{Data $\rightarrow$ Data Analysis}. Select \textbf{Regression}.
\vspace{7pt}

Use desktop Excel: the standard ToolPak Regression command is unavailable in Excel for the web.
\end{frame}

\begin{frame}{Step 6. Fill in the Regression dialog}
\begin{center}
\begin{tabular}{ll}
\toprule
Field or option & What to enter or select\\
\midrule
Input Y Range & \xl{Data!$B$1:$B$101}\\
Input X Range & \xl{Data!$C$1:$C$101}\\
Labels & Check: our selected ranges include row 1\\
Constant is Zero & \textbf{Leave unchecked}\\
Confidence Level & 95\%\\
New Worksheet Ply & Select; enter \xl{OLS_simple}\\
Residuals & Check\\
Residual Plots & Check\\
\bottomrule
\end{tabular}
\end{center}
Click \textbf{OK}. Excel creates a new sheet containing the regression results.
\vspace{5pt}

\textbf{Do not include school ID or income in this first X range.}
\end{frame}

\begin{frame}{Step 7. Read the regression summary}
\begin{columns}[T]
\begin{column}{.40\textwidth}
\begin{tabular}{lr}
\toprule
Regression Statistics & Value\\
\midrule
Multiple R & 0.7731\\
R Square & 0.5977\\
Adjusted R Square & 0.5936\\
Standard Error & 13.3138\\
Observations & 100\\
\bottomrule
\end{tabular}
\end{column}
\begin{column}{.56\textwidth}
\begin{itemize}
\setlength{\itemsep}{7pt}
\item \textbf{R Square}: the fitted line accounts for about 59.77\% of the sample variation in scores.
\item \textbf{Standard Error}: the SER, measured in score points. It is not the standard error of the slope.
\item \textbf{Multiple R}: here, $\sqrt{R^2}$. It is nonnegative, even though the slope is negative.
\item \textbf{Observations}: check that Excel used all 100 schools.
\end{itemize}
\end{column}
\end{columns}
\vspace{8pt}
The ANOVA block decomposes variation. For this exercise, focus first on the summary and coefficient rows.
\end{frame}

\begin{frame}{Step 8. Find the coefficients and their uncertainty}
\begin{center}
\small
\begin{tabular}{lrrrrrr}
\toprule
 & Coefficient & Std. Error & t Stat & P-value & Lower 95\% & Upper 95\%\\
\midrule
Intercept & 749.0587 & 8.0760 & 92.7514 & $2.75\!\times\!10^{-97}$ & 733.0322 & 765.0853\\
X\_class\_size & $-4.4236$ & 0.3666 & $-12.0667$ & $4.33\!\times\!10^{-21}$ & $-5.1511$ & $-3.6961$\\
\bottomrule
\end{tabular}
\end{center}
\begin{itemize}
\setlength{\itemsep}{8pt}
\item Read the \textbf{X\_class\_size} row for the estimated slope and its standard error.
\item Excel's \textbf{P-value} tests a coefficient equal to zero against a two-sided alternative.
\item Scientific notation such as \xl{4.33E-21} means $4.33\times10^{-21}$, not zero.
\item These are conventional, homoskedastic standard errors. ToolPak does not automatically report heteroskedasticity-robust standard errors.
\end{itemize}
\end{frame}

\begin{frame}{Translate the estimated line into a sentence}
\[
\widehat{\text{score}}=749.0587-4.4236\,\text{class size}.
\]
\begin{itemize}
\setlength{\itemsep}{9pt}
\item A school with an average class size larger by one student is predicted to score about \textbf{4.42 points lower}.
\item A difference of five students corresponds to $5(-4.4236)=-22.12$ predicted score points.
\item The intercept is the predicted score at a class size of zero. That is outside our data range, so it has no useful school interpretation here.
\item The fitted association does not show what would happen if a school changed its class size. Family background may differ across schools.
\end{itemize}
\end{frame}

\begin{frame}{Step 9. Reproduce the coefficients with Excel formulas}
Create a new sheet named \xl{Calculations}. Enter these labels and formulas:
\begin{center}
\begin{tabular}{lll}
\toprule
Cell in column A & Formula in column B & Expected value\\
\midrule
A2: Intercept & \xl{=INTERCEPT(Data!B2:B101,Data!C2:C101)} & 749.0587\\
A3: Slope & \xl{=SLOPE(Data!B2:B101,Data!C2:C101)} & $-4.4236$\\
A4: n & \xl{=COUNT(Data!B2:B101)} & 100\\
A5: R squared & \xl{=RSQ(Data!B2:B101,Data!C2:C101)} & 0.5977\\
\bottomrule
\end{tabular}
\end{center}
\begin{itemize}
\setlength{\itemsep}{7pt}
\item In each function, enter the $Y$ range first and the $X$ range second.
\item Show four decimals for readability, but keep the full underlying values.
\item Unlike ToolPak output, these formulas update if the data change. Rerun ToolPak after changing data.
\end{itemize}
\end{frame}

\begin{frame}{Step 10. Calculate fitted values and residuals}
Return to \xl{Data}. Add three column labels and enter formulas in row 2:
\begin{center}
\begin{tabular}{ll}
\toprule
New column & Formula in the first data row\\
\midrule
E1: Y\_fitted & E2: \xl{=Calculations!$B$2+Calculations!$B$3*C2}\\
F1: Residual & F2: \xl{=B2-E2}\\
G1: Residual\_squared & G2: \xl{=F2^2}\\
\bottomrule
\end{tabular}
\end{center}
\begin{enumerate}
\setlength{\itemsep}{6pt}
\item Select E2:G101 and use \textbf{Home $\rightarrow$ Fill $\rightarrow$ Down}.
\item The \xl{$} signs keep the coefficient cells fixed; \xl{C2} changes to \xl{C3}, and so on.
\item Check school 1: $\hat y_1\approx626.0816$ and $\hat u_1\approx-15.2630$.
\end{enumerate}
\textbf{Meaning:} school 1 scored about 15.26 points below its fitted value.
\end{frame}

\begin{frame}{Step 11. Inspect the residual plot}
\centering
\begin{tikzpicture}
\begin{axis}[schoolplot,ylabel={Residual (score points)},ymin=-40,ymax=40]
\addplot[only marks,mark=*,mark size=1.5pt,blue] table[x=X_class_size,y expr=\thisrow{Y_score}-(749.0587322002461-4.423638908332379*\thisrow{X_class_size}),col sep=comma]{\datapath};
\addplot[red,thick,domain=16:28]{0};
\end{axis}
\end{tikzpicture}
\par\small Use X values C2:C101 and Y values F2:F101 in an XY scatterplot. Look for curvature, unusually large residuals, or changing spread; this plot alone cannot establish causality.
\end{frame}

\begin{frame}{Step 12. Calculate SSR, RMSE and SER}
On \xl{Calculations}, continue the table:
\begin{center}
\begin{tabular}{lll}
\toprule
Label & Formula & Expected value\\
\midrule
A6: SSR & B6: \xl{=SUM(Data!G2:G101)} & 17371.2137\\
A7: RMSE & B7: \xl{=SQRT(B6/B4)} & 13.1800\\
A8: SER & B8: \xl{=SQRT(B6/(B4-2))} & 13.3138\\
\bottomrule
\end{tabular}
\end{center}
\[
SSR=\sum_{i=1}^{100}\hat u_i^2,\text{ and }
RMSE=\sqrt{\frac{SSR}{100}}.
\]
\[
SER=\sqrt{\frac{SSR}{100-2}}.
\]
Two estimated coefficients leave $100-2=98$ residual degrees of freedom. Excel's summary \textbf{Standard Error} is SER, not RMSE.
\end{frame}

\begin{frame}{Step 13. Test whether the slope is zero}
\[
H_0:\beta_1=0,\text{ and }H_1:\beta_1\ne0.
\]
On \xl{Calculations}, enter:
\begin{center}
\begin{tabular}{llr}
\toprule
Label & Formula & Value\\
\midrule
A9: SE of slope & B9: \xl{=B8/SQRT(DEVSQ(Data!C2:C101))} & 0.3666\\
A10: t statistic & B10: \xl{=B3/B9} & $-12.0667$\\
\bottomrule
\end{tabular}
\end{center}
\[
t_{\mathrm{obs}}=\frac{-4.4236-0}{0.3666}\approx-12.07.
\]
\begin{itemize}
\item Using the normal approximation already covered: $|t_{\mathrm{obs}}|>1.96$, so reject $H_0$ at 5\%.
\item Excel reports a two-sided p-value of about $4.33\times10^{-21}$ using a $t$ reference distribution with 98 degrees of freedom. The decision is the same.
\end{itemize}
\end{frame}

\begin{frame}{Step 14. Build and interpret a confidence interval}
\begin{center}
\begin{tabular}{lll}
\toprule
Label & Formula & Value\\
\midrule
A11: Normal critical value & B11: \xl{=1.96} & 1.96\\
A12: Lower endpoint & B12: \xl{=B3-B11*B9} & $-5.1422$\\
A13: Upper endpoint & B13: \xl{=B3+B11*B9} & $-3.7051$\\
\bottomrule
\end{tabular}
\end{center}
\[
\hat\beta_1\pm1.96\,SE(\hat\beta_1)\approx[-5.14,-3.71].
\]
\begin{itemize}
\setlength{\itemsep}{7pt}
\item $1.96$ leaves 95\% of a standard normal distribution between $-1.96$ and $1.96$.
\item Excel's \textbf{Lower 95\%} and \textbf{Upper 95\%} use $1.9845$ from its $t$ reference distribution: $[-5.1511,-3.6961]$.
\item Both intervals exclude zero. Excluding zero is equivalent to rejecting the corresponding two-sided 5\% test.
\end{itemize}
\end{frame}

\begin{frame}{Write a short regression report}
\begin{block}{Example}
Across 100 synthetic schools, average class size is negatively associated with average test scores. An additional student per class is associated with approximately 4.42 fewer score points (standard error 0.37; $p<0.001$). The fitted line accounts for 59.77\% of the sample variation in scores, and the SER is 13.31 points.
\end{block}
\begin{itemize}
\setlength{\itemsep}{8pt}
\item Include the sample size, units, coefficient and uncertainty.
\item Distinguish statistical significance from the size of the association.
\item State the limitation: schools differ in family background and other factors, so this regression alone does not identify a causal class-size effect.
\end{itemize}
\end{frame}

\section{Adding Z}
\begin{frame}{Now introduce $Z$: family income}
\begin{itemize}
\setlength{\itemsep}{8pt}
\item Column D contains a new explanatory variable for the same 100 schools: \xl{Z_income_kEUR}.
\item A value of 40 means an average annual family income of \euro{}40,000.
\item Income may affect learning resources and therefore test scores. It may also be related to class size.
\item Every income value must remain matched to the same school ID, score and class size.
\end{itemize}
\begin{block}{Check it in Excel}
\xl{=AVERAGE(Data!D2:D101)} gives 36.6975.\\[4pt]
\xl{=CORREL(Data!C2:C101,Data!D2:D101)} gives $-0.6976$.
\end{block}
The dataset already includes $Z$, so the transition to Lecture 5 uses exactly the same observations.
\end{frame}

\begin{frame}{Schools with smaller classes also have higher family income}
\centering
\begin{tikzpicture}
\begin{axis}[schoolplot,ylabel={Family income (thousand euros/year)},ymin=15,ymax=60]
\addplot[only marks,mark=*,mark size=1.5pt,blue] table[x=X_class_size,y=Z_income_kEUR,col sep=comma]{\datapath};
\end{axis}
\end{tikzpicture}
\par\small Recreate this graph with C2:C101 on the horizontal axis and D2:D101 on the vertical axis. Sample correlation: $-0.6976$.
\end{frame}

\begin{frame}{Why could omitting income change the slope?}
\begin{enumerate}
\setlength{\itemsep}{9pt}
\item Income is a determinant of scores in our simulated example.
\item Income is negatively correlated with class size.
\item A comparison of large and small classes therefore also compares schools with different family incomes.
\end{enumerate}
\begin{block}{Connection to Lecture 5: omitted-variable bias}
Both conditions matter: $Z$ helps determine $Y$, and $Z$ is correlated with $X$. The simple slope combines the class-size relationship with the income differences associated with class size.
\end{block}
With a positive income effect and a negative income--class-size relationship, we expect the simple slope to be more negative than the slope controlling for income.
\end{frame}

\begin{frame}{Prepare the multiple regression for Lecture 5}
\[
Y_i=\alpha_0+\alpha_1X_i+\alpha_2Z_i+\varepsilon_i.
\]
\begin{enumerate}
\setlength{\itemsep}{7pt}
\item Return to \textbf{Data $\rightarrow$ Data Analysis $\rightarrow$ Regression}.
\item Keep Input Y Range: \xl{Data!$B$1:$B$101}.
\item Change Input X Range to \xl{Data!$C$1:$D$101}: two adjacent columns, class size \emph{and} income.
\item Check \textbf{Labels}, keep the intercept, and use a new output sheet named \xl{OLS_multiple}.
\end{enumerate}
\begin{block}{The new question}
How are class size and scores related \textbf{holding average family income constant}? We use different symbols, $\alpha$, because these coefficients belong to a different model.
\end{block}
\end{frame}

\begin{frame}{Preview: what changes when we include income?}
\begin{center}
\begin{tabular}{lrr}
\toprule
 & Simple regression & Multiple regression\\
\midrule
Intercept & 749.0587 & 646.2105\\
Class-size coefficient & $-4.4236$ & $-2.1867$\\
Income coefficient & Not included & 1.4781\\
$R^2$ & 0.5977 & 0.7590\\
SER & 13.3138 & 10.3588\\
Observations & 100 & 100\\
\bottomrule
\end{tabular}
\end{center}
\begin{itemize}
\setlength{\itemsep}{7pt}
\item Holding income constant, one additional student per class is associated with about 2.19 fewer score points.
\item Holding class size constant, an additional \euro{}1,000 of income is associated with about 1.48 more score points.
\item Adding a regressor cannot lower in-sample $R^2$. A higher $R^2$ alone does not establish causality.
\end{itemize}
\end{frame}

\section{Practice}
\begin{frame}{Your turn: six checks before leaving Excel}
\begin{enumerate}
\setlength{\itemsep}{9pt}
\item Predict the score for a school with an average class size of 20 using the simple regression.
\item What happens to the fitted score if class size rises from 20 to 25?
\item Calculate \xl{=SUM(Data!F2:F101)}. Why is it approximately zero?
\item Which number is the slope's standard error: 13.3138 or 0.3666?
\item Why does the normal-approximation confidence interval exclude zero?
\item Compare the two class-size coefficients. What does adding income change about the comparison?
\end{enumerate}
\end{frame}

\begin{frame}{Answers and a troubleshooting checklist}
\begin{columns}[T]
\begin{column}{.50\textwidth}
\begin{block}{Answers}
\begin{enumerate}
\setlength{\itemsep}{5pt}
\item $749.0587-4.4236(20)\approx660.5860$.
\item About 22.1182 points lower.
\item OLS with an intercept makes residuals sum to zero, up to rounding.
\item 0.3666; 13.3138 is the SER.
\item The estimate is more than 1.96 standard errors from zero.
\item We now compare class sizes holding income constant.
\end{enumerate}
\end{block}
\end{column}
\begin{column}{.46\textwidth}
\begin{block}{If your output differs}
\begin{itemize}
\setlength{\itemsep}{7pt}
\item Confirm Y is column B and X is column C.
\item Use rows 1--101 with \textbf{Labels} checked, or 2--101 without it.
\item Leave \textbf{Constant is Zero} unchecked.
\item Exclude IDs and missing cells.
\item Keep full precision; rerun ToolPak after changing data.
\end{itemize}
\end{block}
\end{column}
\end{columns}
\end{frame}

\begin{frame}{Behind the data: how the 100 schools were created}
\small
Independent draws across schools, with fixed seed 251174:
\begin{align*}
X_i&\sim\operatorname{Uniform}(16,28),\\
V_i&\sim\operatorname{Uniform}(-10,10),\text{ and }\varepsilon_i\sim N(0,10^2),\\
Z_i&=70-1.5X_i+V_i,\\
U_i&=1.5(Z_i-40)+\varepsilon_i,\\
Y_i&=700-2X_i+U_i=640-2X_i+1.5Z_i+\varepsilon_i.
\end{align*}
\begin{itemize}
\item $X_i$, $V_i$ and $\varepsilon_i$ are drawn independently. Stored class sizes have one decimal, incomes two, and scores four; see the generation script for the rounding order.
\item This has the form $Y=b_0+b_1X+U$, with $b_0=700$ and $b_1=-2$; the omitted income component makes $E(U\mid X)\ne0$.
\item Therefore $-2$ is the class-size effect holding income constant in this simulation, not the slope of the simple conditional mean. Regressing on $X$ alone need not recover it.
\end{itemize}
\end{frame}

\begin{frame}{Files and Excel help}
\begin{itemize}
\setlength{\itemsep}{9pt}
\item \xl{data/lecture4v3/school_regression.xlsx}: fixed data and a worked formula sheet.
\item \xl{data/lecture4v3/school_data.csv}: the same 100 observations in a portable format.
\item \xl{data/lecture4v3/generate_data.py}: reproduces the synthetic observations and independently checks the regression results.
\item Microsoft Support: \href{https://support.microsoft.com/en-us/excel/load-the-analysis-toolpak-in-excel}{\textcolor{blue}{Load the Analysis ToolPak in Excel}}.
\item Microsoft Support: \href{https://support.microsoft.com/en-us/excel/use-the-analysis-toolpak-to-perform-complex-data-analysis}{\textcolor{blue}{Use the Analysis ToolPak to perform complex data analysis}}.
\item Microsoft Support: \href{https://support.microsoft.com/en-us/excel/perform-a-regression-analysis}{\textcolor{blue}{Perform a regression analysis}}.
\end{itemize}
\vspace{6pt}
The figures and benchmark results use the saved dataset. Menu names can vary by Excel version and language.
\end{frame}
\end{document}
